%!TEX root = ../main.tex
%----------------------------------------------------------------------
%  Section 5: Data and Baseline Models
%  Paper:   Theory-Informed Kernel Ridge Regression for Asset Pricing
%  Authors: Amedeo Andriollo and Juan F. Imbet
%----------------------------------------------------------------------

\section{Data and Baseline Models}\label{sec:empirical}

%======================================================================
\subsection{Data}\label{sub:data}
%======================================================================

Our primary dataset consists of monthly stock returns from the Center for Research in Security Prices (CRSP) monthly stock file.
The sample includes ordinary common shares (share codes 10 and 11) listed on NYSE, AMEX, and NASDAQ (exchange codes 1, 2, and 3).
Excess returns are computed by subtracting the one-month Treasury bill rate from the Fama-French data library.
The effective sample period begins in July~1963---when both CRSP and Compustat coverage become reliable---and ends in December~2023, spanning 726~months and approximately 18,000 unique securities.
This start date matches \citet{gu2020empirical} and the bulk of the cross-sectional return prediction literature.

We pair returns with 220 firm-level predictive characteristics constructed from CRSP returns, Compustat fundamentals, and the Open Source Asset Pricing (OSAP) dataset of \citet{jensen2023machine}.
The OSAP dataset provides a standardized, publicly available panel of signed characteristics at the security-month level, constructed following the original papers' methodologies and validated against published results.
Supplementary variables include CRSP-derived signals (price, size, short-term reversal) and Compustat-derived fundamentals (investment-to-capital, ROE, leverage, ROA, gross profitability, asset growth, R\&D intensity) that enter the structural restriction penalties directly.
All characteristics are signed so that higher values predict higher expected returns.
This characteristic set subsumes and extends the 94~characteristics used in \citet{gu2020empirical}.

We merge the OSAP characteristics with CRSP returns using the permanent security identifier (permno).
To ensure comparability across variables measured in different units, we cross-sectionally rank each characteristic to the unit interval $[0,1]$ each month.
Ranking is a monotone transformation that preserves the ordinal information in each characteristic while eliminating outliers and ensuring that all variables enter the kernel on a common scale.

The structural restrictions in Section~\ref{sec:restrictions} additionally require firm-level accounting data from Compustat.
We obtain annual fundamentals (525,028 firm-years, 40,678 unique firms, 1950--2025) and quarterly fundamentals (1.87~million firm-quarters, 39,568 unique firms, 1961--2025) from WRDS.
These are merged with CRSP via the CRSP-Compustat linking table using primary link types (LC and LU) with an effective link date filter.
Key accounting variables include total assets, book equity, sales, income before extraordinary items, capital expenditures, property-plant-equipment, depreciation, long-term and short-term debt, cash holdings, R\&D expenditure, and selling-general-administrative expenses.
Our Compustat sample includes both active and inactive firms to avoid survivorship bias.

We augment firm-level characteristics with macroeconomic state variables that serve two roles: they enter the kernel as predictors of expected returns, and they supply inputs to the structural restriction penalties in Section~\ref{sec:restrictions}.
These variables are drawn from four sources.

From the Federal Reserve Economic Data (FRED), we obtain the 3-month Treasury bill rate (TB3MS), the 10-year Treasury yield (GS10), Moody's AAA and BAA corporate bond yields, the Consumer Price Index (CPIAUCSL), industrial production (INDPRO), and the unemployment rate (UNRATE).
The term spread (GS10 minus TB3MS) and the default spread (BAA minus AAA) are computed as differences.
From the National Income and Product Accounts via FRED, we obtain personal consumption expenditures on nondurable goods and services, the PCE price index, and population, which we combine to construct real per capita consumption growth---the central input to the thirteen consumption-based Euler equation restrictions.

The intermediary capital ratio and intermediary capital risk factor of \citet{hekellymanela2017} are obtained from the authors' website at quarterly frequency and interpolated to monthly.
The consumption-wealth ratio $\mathit{cay}$ of \citet{lettauludvigson2001} is constructed following their Dynamic OLS methodology from FRED series on household net worth (TNWBSHNO), employee compensation (COE), and consumption, using 8~leads and lags of the differenced regressors.
The \citet{bakerwurgler2006} investor sentiment index (orthogonalized version) is obtained from the authors' data repository and covers 1965--2023.
The CBOE Volatility Index (VIX) is obtained from FRED and is available from January~1990; we aggregate daily values to monthly averages.
The excess bond premium of \citet{gilchristzakrajsek2012} is proxied by the Moody's BAA spread over the 10-year Treasury (BAAFFM series from FRED).
Additional macro predictors---the dividend-price ratio, earnings-price ratio, aggregate book-to-market, net equity expansion, Treasury bill rate, long-term yield, and stock variance---are obtained from the \citet{welchgoyal2008} dataset, which covers 1871--2024 at monthly frequency.

Several structural restrictions require variables beyond the standard characteristic and macro sets.
For the production-based restrictions (14--23), we construct leverage (total debt over total assets), R\&D intensity (R\&D expenditure over sales), SGA intensity (selling, general, and administrative expenses over total assets), and the collateral ratio (PP\&E over total debt) from Compustat annual fundamentals.
The hiring rate of \citet{belolinbazdresch2014} is obtained from the OSAP characteristic panel.
The dividend yield used in the costly external finance restriction~(21) is the short-term dividend yield from the OSAP panel.

For intermediary restrictions, we supplement the HKM factor with broker-dealer leverage from the Federal Reserve's Financial Accounts (Flow of Funds series BOGZ1FL664090005Q and BOGZ1FL665080003Q), following \citet{adrianetulamuir2014}.
The TED spread (3-month LIBOR minus 3-month Treasury bill, FRED series TEDRATE) proxies funding liquidity in restriction~26; this series ends in January~2022 following the discontinuation of LIBOR.
The betting-against-beta (BAB) factor of \citet{frazzinipedersen2014} is obtained from AQR's data library and covers 1930--2025 at monthly frequency.
Household equity net flows (FRED series HNOCESQ027S, quarterly, forward-filled to monthly) proxy capital flows in restrictions~29 and~40.

For behavioral restrictions, the capital gains overhang of restriction~48 is computed as $(P_t - \mathrm{RP}_t)/P_t$, where the reference price $\mathrm{RP}_t$ is an exponentially weighted moving average of past prices with a 12-month half-life.
The salience measure of restriction~51 is $|R_{i,t} - R_{m,t}|/\sigma_{\mathrm{cross},t}$, the absolute deviation of the stock's return from the market return scaled by the cross-sectional standard deviation.
The past average return of restriction~49 is the trailing 36-month rolling mean of monthly returns.

For macro-finance restrictions, the 10-year breakeven inflation rate (FRED series T10YIEM, available from 2003) proxies the inflation risk premium in restriction~53.
The monetary policy surprise of restriction~54 is the residual from an AR(1) model of monthly changes in the effective federal funds rate (FRED series FEDFUNDS).
The $q$-factor expected growth factor ($R_{\mathrm{EG}}$) of \citet{houxuezhang2015} is obtained from the Macro Finance Society dataset on WRDS.

We require a minimum of 15 non-missing characteristics per stock-month observation.
Observations with fewer non-missing characteristics are dropped.
Returns are winsorized at the 0.1st and 99.9th percentiles each month to limit the influence of extreme observations while preserving the bulk of the return distribution.
After applying all filters, the sample contains 1,778,850 stock-month observations spanning 18,000 unique securities, covering an average of approximately 2,450~stocks per month.

Not all macro state variables are available for the full sample period.
The VIX is available only from 1990, so the two variance and volatility restrictions (44--45) are evaluated on the post-1990 subsample.
The HKM intermediary factors are available quarterly from 1970 to 2012; restrictions~24--27 are evaluated on this subsample.
The Baker-Wurgler sentiment index is available from 1965 to 2023.
The TED spread is available from 1986 to 2022 (discontinued with LIBOR).
The 10-year breakeven inflation rate is available from 2003.
The HXZ expected growth factor covers 1967--2019.
All other macro variables cover the full 1963--2023 sample period.

%======================================================================
\subsection{Baseline Models}\label{sub:baselines}
%======================================================================

We compare Theory-KRR against five baseline models spanning the spectrum from simple linear methods to a flexible nonparametric estimator.
All baselines use the same data, the same rolling-window procedure described in Section~\ref{sub:cv_strategy}, and the same temporal cross-validation for hyperparameter selection.
This ensures that performance differences reflect the models' functional form and regularization, not differences in data or evaluation protocol.

\begin{enumerate}[label=\arabic*.]
\item \textbf{Historical mean.}
The simplest benchmark: $\hat{f}(\mathbf{X}_{i,t}) = \bar{R}$, the pooled cross-sectional average of excess returns in the training sample.
This is the natural benchmark for $R^2_{\mathrm{OOS}}$, which measures improvement over the historical mean.

\item \textbf{OLS.}
Linear regression of next-month excess returns on the full vector of characteristics: $\hat{f}(\mathbf{X}_{i,t}) = \hat{\bm{\beta}}^{\top}\mathbf{X}_{i,t}$.
OLS is prone to overfitting in small training samples but provides a benchmark for the value of regularization and nonlinearity.

\item \textbf{Ridge regression.}
OLS with an $\ell_2$ penalty on the coefficient vector: the objective is $\frac{1}{n}\sum_s(R_s - \bm{\beta}^\top \mathbf{X}_s)^2 + \lambda\|\bm{\beta}\|_2^2$.
Ridge regression shrinks coefficients toward zero without enforcing sparsity.
The regularization parameter $\lambda$ is selected on the validation sample over a grid of 20 log-spaced values.

\item \textbf{Lasso.}
OLS with an $\ell_1$ penalty: the objective is $\frac{1}{n}\sum_s(R_s - \bm{\beta}^\top \mathbf{X}_s)^2 + \lambda\|\bm{\beta}\|_1$.
The $\ell_1$ penalty performs variable selection by setting some coefficients exactly to zero.
The regularization path is computed by coordinate descent \citep{friedman2010regularization} and the optimal $\lambda$ is selected on the validation sample.

\item \textbf{Elastic net.}
A convex combination of $\ell_1$ and $\ell_2$ penalties with mixing parameter $\alpha = 0.5$.
The regularization path and $\lambda$ selection follow the same procedure as the Lasso.

\item \textbf{Standard KRR.}
The Gaussian kernel ridge regression estimator with RKHS-norm regularization only: this is Theory-KRR with all $\mu_j = 0$.
Standard KRR captures nonlinearities and interactions through the Gaussian kernel but receives no guidance from economic theory.
Comparing Theory-KRR to standard KRR isolates the marginal contribution of the structural penalties.
\end{enumerate}

%======================================================================
\subsection{Evaluation Metrics}\label{sub:evaluation}
%======================================================================

We evaluate predictive performance along two dimensions: statistical accuracy and economic significance.
Following Campbell and Thompson~(2008), we measure predictive accuracy by the out-of-sample $R^2$:
\begin{equation}\label{eq:r2oos}
R^2_{\mathrm{OOS}} = 1 - \frac{\displaystyle\sum_{(i,t) \in \mathcal{T}_{\mathrm{OOS}}} \bigl(R_{i,t+1} - \hat{f}(\mathbf{X}_{i,t})\bigr)^2}{\displaystyle\sum_{(i,t) \in \mathcal{T}_{\mathrm{OOS}}} \bigl(R_{i,t+1} - \bar{R}\bigr)^2},
\end{equation}
where $\mathcal{T}_{\mathrm{OOS}}$ indexes the out-of-sample observations and $\bar{R}$ is the historical mean computed from the training sample.
A positive $R^2_{\mathrm{OOS}}$ indicates that the model predicts better than the historical mean.
In the cross-section of individual stock returns, $R^2_{\mathrm{OOS}}$ values above 1\% are economically meaningful given the high idiosyncratic variance of monthly returns.

To assess economic significance, we form decile portfolios each month by sorting stocks on their predicted expected returns $\hat{f}(\mathbf{X}_{i,t})$.
The long-short portfolio goes long the top decile and short the bottom decile, with value-weighted holdings within each decile.
We report annualized Sharpe ratios:
\begin{equation}\label{eq:sharpe}
\mathrm{SR} = \frac{\hat{\mu}_{L-S}}{\hat{\sigma}_{L-S}} \times \sqrt{12},
\end{equation}
where $\hat{\mu}_{L-S}$ and $\hat{\sigma}_{L-S}$ are the sample mean and standard deviation of the monthly long-short return.
Higher Sharpe ratios indicate that the model's predictions generate more profitable trading strategies per unit of risk.

We assess the statistical significance of performance differences using the Diebold-Mariano test, which compares squared forecast errors between Theory-KRR and each baseline with Newey-West standard errors (6 lags) to account for serial correlation.
All reported standard errors use the Newey-West HAC estimator with 6 lags unless otherwise noted.

%======================================================================
\subsection{Summary Statistics}\label{sub:summary}
%======================================================================

Table~\ref{tab:summary} documents the key features of our dataset.
Panel~A shows that the average monthly excess return is 2.0\% with a standard deviation of 14.6\%, consistent with well-known properties of the cross-section.
All firm-level characteristics are cross-sectionally ranked to the unit interval each month, producing uniform marginal distributions with mean 0.50 and standard deviation 0.29.
Panel~B reports that macroeconomic state variables display high persistence: the intermediary capital ratio of \citet{hekellymanela2017} has an autocorrelation coefficient of 0.98, the consumption-wealth ratio $cay$ has an autocorrelation of 0.99, and the Baker-Wurgler sentiment index has an autocorrelation of 0.99.
Consumption growth is the exception, with negative autocorrelation ($-0.47$), reflecting its near-i.i.d.\ behavior at monthly frequency.
Panel~C confirms that sample coverage peaks in the 1990s at approximately 3,400 stocks per month and declines to approximately 2,200 in the 2010s, reflecting the secular decline in the number of publicly listed U.S.\ firms.
